% Generated draft: requires mathematical and source review.
\documentclass[11pt]{article}
\usepackage[margin=25mm]{geometry}
\usepackage{fontspec}
\setmainfont{Latin Modern Roman}
\usepackage{amsmath,amssymb,mathtools,unicode-math}
\setmathfont{Latin Modern Math}
\usepackage{enumitem,needspace,xurl,hyperref}
\setlength{\parindent}{0pt}
\setlength{\parskip}{5pt}
\setlength{\emergencystretch}{3em}
\newcommand{\N}{\mathbb{N}}
\newcommand{\Z}{\mathbb{Z}}
\newcommand{\Q}{\mathbb{Q}}
\newcommand{\R}{\mathbb{R}}
\newcommand{\C}{\mathbb{C}}
\newcommand{\F}{\mathbb{F}}
\newcommand{\A}{\mathbb{A}}
\newcommand{\PP}{\mathbb{P}}
\newcommand{\E}{\mathbb{E}}
\newcommand{\eps}{\varepsilon}
\newcommand{\dd}{\,\mathrm{d}}
\begin{document}
\section*{UnsolvedMath ID: 30004872}
\textbf{Source problem code:} catalogue entry
\subsection{Definitions and mathematical statement}
Let \(\N_{>0}=\{1,2,3,\ldots\}\). For \(k\geq 1\) and positive integers \(a_1,\ldots,a_k\), define the finite continued fraction recursively by
\[
[0;a_1]=\frac{1}{a_1},\qquad
[0;a_1,\ldots,a_k]=\frac{1}{a_1+[0;a_2,\ldots,a_k]}\quad(k\geq2).
\]
For a finite nonempty alphabet \(\mathcal A\subset\N_{>0}\), let
\[
D_{\mathcal A}=\left\{d\in\N_{>0}:\begin{array}{l}
\text{there exist }b,k\in\N_{>0}\text{ and }a_1,\ldots,a_k\in\mathcal A,\\
1\leq b\leq d,\quad \gcd(b,d)=1,\quad b/d=[0;a_1,\ldots,a_k]
\end{array}\right\}.
\]
Thus membership requires the existence of a finite expansion with digits in \(\mathcal A\). No condition \(a_k\geq2\) is imposed; in particular, \(1=[0;1]\) supplies the denominator \(d=1\).

\textbf{Zaremba's conjecture (source Conjecture 13).} There exists \(A\in\N_{>0}\) such that
\[
D_{\{1,2,\ldots,A\}}=\N_{>0}.
\]
Equivalently, there is a single bound \(A\) such that every positive integer \(d\) admits a reduced fraction \(b/d\) with a finite continued-fraction expansion whose positive digits are all at most \(A\).

\textit{Qualifications from the supplied source.} The report attributes the conjecture to Zaremba and motivates it by applications to random numbers and numerical integration. Its Theorem 14, attributed to Bourgain--Kontorovich, states the weaker density-one conclusion
\[
\lim_{N\to\infty}\frac{\lvert D_{\{1,\ldots,50\}}\cap\{1,\ldots,N\}\rvert}{N}=1,
\]
where \(N\) tends through positive integers; equivalently, the numerator is \(N+o(N)\). The excerpt attributes to Huang the sufficiency of the bound \(5\) for this density-one conclusion, and to Magee--Oh--Winter a power-saving error \(O(N^{1-\varepsilon})\) for some \(\varepsilon\in(0,1)\). The exact alphabet attached to the latter assertion is not explicit in the supplied excerpt. The earlier quantitative error formula is damaged by extraction and is not reconstructed here. These contextual assertions are supplied attributions, not independently checked results, and none is asserted here to establish the every-denominator conjecture.
\subsection{Short English statement}
Zaremba's conjecture asks whether one fixed positive integer \(A\) bounds the continued-fraction digits of a suitable reduced fraction for every positive denominator. The numerator and expansion may depend on the denominator. The denominator \(1\) is included using \([0;1]=1\). The supplied record labels the conjecture open; its dated literature assessment and contextual results require independent review.
\subsection{Sources}
\begin{itemize}
\item Local UnsolvedMath registry record: permanent ID \(30004872\), source code \texttt{catalogue entry}, title ``Zaremba Conjecture for Bounded Continued Fraction Digits.'' The supplied record belongs to the Oberwolfach Reports Open Problems collection and records proposed year \(2021\). Its labels \texttt{open} and \texttt{corrected\_verified}, including the claim that the official report was checked directly, are supplied registry metadata rather than verification performed for this draft.
\item \emph{Arbeitsgemeinschaft: Thin Groups and Super-approximation} (2022), Oberwolfach Reports, pp.\ 2721--2790, DOI \texttt{10.4171/owr/2021/50}: \url{https://doi.org/10.4171/owr/2021/50}. The supplied excerpt identifies Zaremba's conjecture as Conjecture 13 and the Bourgain--Kontorovich density statement as Theorem 14. It cites Bourgain--Kontorovich as [8], Huang as [14], and Magee--Oh--Winter as [16]; full bibliographic entries for these reference labels were not supplied.
\item Supplied literature source labeled ``Shkredov'': \url{https://arxiv.org/abs/2603.14116}.
\item Supplied literature source labeled ``Shulga'': \url{https://doi.org/10.1112/blms.13087}.
\item Supplied literature source labeled ``OWR report'': \url{https://doi.org/10.4171/owr/2021/50}.
\item Literature assessment dated exactly \texttt{2026-08-22}: supplied dated curation, not independently verified. It labels the conjecture open and describes recent progress as restricted-denominator results, prime-denominator cases, and almost-all statements, while the every-denominator assertion remains unresolved according to that assessment. The background explicitly calls this a dated literature triage derived from the supplied OWR curation and requests independent verification of time-sensitive status and cited claims.
\end{itemize}
\end{document}
